\documentclass[11pt]{article}
\usepackage[T1]{fontenc}
\usepackage[utf8]{inputenc}
\usepackage{amsmath,amssymb,amsthm}
\usepackage[margin=1in]{geometry}
\usepackage[hidelinks]{hyperref}
\hypersetup{
  pdftitle={A Four-Dimensional Witness for the Modular-Balance Gap},
  pdfauthor={Xiaojian Zeng}
}

\newtheorem{theorem}{Theorem}

\newcommand{\R}{\mathbb{R}}
\newcommand{\h}{\mathfrak{h}}
\newcommand{\q}{\mathfrak{q}}
\newcommand{\ad}{\operatorname{ad}}
\newcommand{\tr}{\operatorname{tr}}

\title{A Four-Dimensional Witness for the Modular-Balance Gap}
\author{Xiaojian Zeng\\
PBC School of Finance, Tsinghua University, 100190\\
\texttt{xiaojianzeng@foxmail.com}}
\date{}

\begin{document}
\maketitle

\begin{abstract}
We give an exact four-dimensional real Lie algebra with a bracket-generating
hyperplane (V) that admits a complement \(q\) satisfying the vertical
trace condition \(\kappa_{\q}|_V=0\), although no complement of the same
\(V\) is modularly balanced.  Every complement is parametrized by three real
numbers, and the two balance conditions reduce respectively to
\(x=y=0\) and \(1+x+y=0\).  Thus the construction proves the bounded
dimension-at-most-four existence assertion and answers the existence clause
of Question~1 of Nicolussi Golo, Pinamonti, and Warhurst affirmatively.  It
does not address the question about normal harmonic maps, classify such
examples, or assert minimality of the dimension.
\end{abstract}

\section{Context and statement}

The classical theory of harmonic morphisms includes the foundational
characterizations of Fuglede and Ishihara~\cite{Fuglede1978,Ishihara1979}.  In the
sub-Riemannian Lie-group setting, Nicolussi Golo, Pinamonti, and
Warhurst introduced a modular mismatch attached to a polarized Lie algebra
and a vector-space complement~\cite{NicolussiGoloPinamontiWarhurst2026}.
We recall only the algebraic definitions needed here.

Let \(\h\) be a finite-dimensional real Lie algebra, let \(V\leq\h\) be a
linear subspace that Lie-generates \(\h\), and fix a vector-space splitting
\(\h=V\oplus\q\).  Write \(\pi_V\) and \(\pi_{\q}\) for the associated
projections and use \(\ad_B(X)=[B,X]\).  The two projected trace characters
are
\begin{equation}\label{eq:characters}
 \kappa_V^{\q}(B)=\tr\!\left(\pi_V\ad_B|_V\right),
 \qquad
 \kappa_{\q}(B)=\tr\!\left(\pi_{\q}\ad_B|_{\q}\right).
\end{equation}
The complement \(\q\) is modularly balanced precisely when
\begin{equation}\label{eq:balance}
 \kappa_V^{\q}|_{\q}=0
 \qquad\text{and}\qquad
 \kappa_{\q}|_V=0.
\end{equation}
These are Definition~8.2 and equation~(60) of
\cite{NicolussiGoloPinamontiWarhurst2026}.  Theorems~10.5 and~10.7 there
isolate a gap between the simultaneous condition \eqref{eq:balance} and the
necessary vertical condition \(\kappa_{\q}|_V=0\).  Their Question~1 asks
whether the latter can hold for some complement while no complement is
modularly balanced.

\begin{theorem}\label{thm:main}
Let \(\h\) have ordered real basis \((e_1,e_2,r,z)\), with the only nonzero
brackets, up to skew-symmetry,
\begin{equation}\label{eq:brackets}
 [e_1,e_2]=e_1+z,
 \qquad [e_1,r]=-r,
 \qquad [z,r]=r.
\end{equation}
Then \eqref{eq:brackets} defines a Lie algebra.  For
\[
 V=\operatorname{span}(e_1,e_2,r)
 \quad\text{and}\quad
 \q_0=\R z,
\]
the subspace \(V\) Lie-generates \(\h\), and
\(\kappa_{\q_0}|_V=0\).  Nevertheless, no vector-space complement of this
same \(V\) is modularly balanced.
\end{theorem}

In particular, Theorem~\ref{thm:main} supplies a dimension-four witness for
the existence clause of Question~1 in
\cite{NicolussiGoloPinamontiWarhurst2026}.

\section{Proof}

First we verify that \eqref{eq:brackets} is a Lie bracket.  For a
skew-symmetric bilinear bracket the Jacobiator is alternating, so it is
enough to check the four triples of distinct basis vectors.  The possibly
nonzero terms reduce to
\begin{align*}
 J(e_1,e_2,r)&=[r,e_1+z]=r-r=0,\\
 J(e_1,e_2,z)&=[z,e_1+z]=0,\\
 J(e_1,r,z)&=[e_1,-r]+[z,-r]=r-r=0,\\
 J(e_2,r,z)&=0.
\end{align*}
Thus Jacobi holds.  Moreover,
\[
 z=[e_1,e_2]-e_1,
\]
so \(V\) Lie-generates all of \(\h\).  The displayed decomposition
\(\h=V\oplus\R z\) is direct.  For every \(B\in V\), the bracket
\([B,z]\) lies in \(V\).  Hence the endomorphism induced by \(\ad_B\) on
the one-dimensional complement has trace zero, proving
\(\kappa_{\q_0}|_V=0\).

It remains to rule out every balanced complement.  Since \(V\) is a
hyperplane, every complement is a line with a unique normalized generator
\begin{equation}\label{eq:w}
 w=z-p,
 \qquad p=xe_1+ye_2+tr,
 \qquad (x,y,t)\in\R^3.
\end{equation}
Denote this complement by \(\q_p=\R w\).  This is the codimension-one form
of the graph chart for complements used in Lemma~10.6 of
\cite{NicolussiGoloPinamontiWarhurst2026}; directly, normalization is
possible and unique because any vector outside \(V\) has nonzero
\(z\)-coordinate.

On the one-dimensional space \(\q_p\), the projected trace of \(\ad_B\)
is the coefficient of \(w\) in \([B,w]\).  Since \(w\) has
\(z\)-coordinate one, this coefficient is also the \(z\)-coordinate in the
original basis.  Expanding \eqref{eq:brackets} gives
\begin{equation}\label{eq:vertical}
 \bigl(\kappa_{\q_p}(e_1),
       \kappa_{\q_p}(e_2),
       \kappa_{\q_p}(r)\bigr)=(-y,x,0).
\end{equation}
Consequently the vertical half of \eqref{eq:balance} holds exactly when
\begin{equation}\label{eq:xyzero}
 x=y=0,
\end{equation}
with \(t\) unrestricted.

For the horizontal half, project \(\ad_w|_V\) onto \(V\) along
\(\q_p\).  In the ordered basis \((e_1,e_2,r)\), its matrix is
\begin{equation}\label{eq:matrix}
 \begin{pmatrix}
 y(1+x) & -x(1+x) & 0\\
 y^2 & -xy & 0\\
 t(y-1) & -xt & 1+x
 \end{pmatrix}.
\end{equation}
Therefore
\begin{equation}\label{eq:horizontal}
 \kappa_V^{\q_p}(w)
 =\tr\!\left(\pi_V\ad_w|_V\right)
 =1+x+y.
\end{equation}
If \(\q_p\) were modularly balanced, \eqref{eq:vertical} would force
\eqref{eq:xyzero}, while \eqref{eq:horizontal} would then equal \(1\), a
contradiction.  The parametrization \eqref{eq:w} covers every complement,
so this proves Theorem~\ref{thm:main}. \qed

\section{Exact verification and scope}

An accompanying Python script recomputes the bracket tensor with exact
SymPy arithmetic~\cite{MeurerEtAl2017}.  It checks all \(4^3=64\) ordered
basis triples, verifies that \(V\) together with \([e_1,e_2]\) has rank
four, derives \eqref{eq:vertical} and \eqref{eq:horizontal}, and reduces
\(1+x+y\) modulo the vertical balance ideal \(\langle x,y\rangle\) to the
nonzero remainder \(1\).  No floating-point approximation or randomized
search enters these checks.  Exact rational certificates have likewise
been used to separate a proposed matrix bound from an attained value without
relying on decimal evidence~\cite{ZengLiuRatnavelu2026}; here the
computer calculation is only a reproducible cross-check of the displayed
proof.

The result is existential.  It does not show that four is the least possible
dimension, classify lower-dimensional polarizations, or address the second
part of Question~1 concerning normal harmonic maps.  The literature
comparison made here is limited to the cited source question and its nearest
stated results; no broader priority claim is made.

\paragraph{Computational and AI assistance.}
OpenAI Codex assisted with the construction search, symbolic derivations,
literature retrieval, verification, and drafting.  SymPy~1.14.0 was used for
the exact cross-check described above.  Neither tool is invoked as a source
of mathematical authority: the complete argument establishing the claim is
given in the proof.

\bibliographystyle{abbrv}
\bibliography{references}

\end{document}
